\section{Quasitraces and proper infiniteness}
\label{sec:preliminaries}

All \(C^*\)-algebras in this paper are complex, and all matrix
algebras carry their usual \(C^*\)-norms. For a unital \(C^*\)-algebra
\(C\), a \emph{\(1\)-quasitrace} is a function
\(\tau:C\to\mathbb C\) satisfying
\begin{enumerate}[label=\textup{(\roman*)}]
\item \(\tau(x^*x)=\tau(xx^*)\ge0\) for every \(x\in C\);
\item \(\tau(h+ik)=\tau(h)+i\tau(k)\) for self-adjoint \(h,k\in C\);
\item the restriction of \(\tau\) to every abelian \(C^*\)-subalgebra
of \(C\) is complex linear.
\end{enumerate}
It is a \emph{\(2\)-quasitrace} if a \(1\)-quasitrace on \(M_2(C)\)
restricts to \(\tau\) on the upper-left corner. It is
\emph{normalized} when \(\tau(1_C)=1\).
The extension then has value two at \(1_{M_2(C)}\): the two
diagonal corner units have the same value by \textup{(i)} and
their values add by \textup{(iii)}.

A unital \(C^*\)-algebra \(C\) is \emph{properly infinite} when it
contains isometries \(s_1,s_2\) with orthogonal ranges. Equivalently,
\[
 s_i^*s_j=\delta_{ij}1_C,\qquad i,j\in\{1,2\}.
\]
It is \emph{stably properly infinite} if \(M_m(C)\) is properly
infinite for some positive integer \(m\). Absence of stable proper
infiniteness is the condition needed below; no assertion about
finiteness of every projection is required.

\begin{theorem}[Quasitrace existence]
\label{thm:existence}
A unital \(C^*\)-algebra \(C\) admits a normalized \(2\)-quasitrace
if and only if \(M_m(C)\) is not properly infinite for every
positive integer \(m\).
\end{theorem}

We use Theorem~\ref{thm:existence} in the formulation of
\cite[Theorem~2.3]{MR}, with the normalization and matrix-extension
conventions fixed at the start of its Section~2. Cuntz's work on
dimension functions \cite{Cuntz1978} precedes the general criterion.
The correspondence between \(2\)-quasitraces and lower semicontinuous
dimension functions is due to Blackadar--Handelman
\cite[Theorems~II.2.2 and II.3.1]{BH}; the extension of
semigroup states underlying this formulation is discussed in
\cite{BR}. In particular, Theorem~\ref{thm:existence} has no exactness,
simplicity, or representation assumption.

For finite-rank complex vector bundles, Hermitian metrics are fixed.
Adjoints, positive endomorphisms, and their square roots are taken
fiberwise. A positive lower bound on the adjoint square of a
continuous bundle map makes it injective with a continuous image
subbundle: if \(W^*W\) is invertible, then
\[
 Q=W(W^*W)^{-1/2},\qquad Q^*Q=I,
\]
has the same image, and \(QQ^*\) is its continuous orthogonal
projection. We will use this observation when passing from
polynomial operators to the Chern-class obstruction.
